% Part III: The JuGeo Pipeline
% Sections 6–9 of the seminal paper (pp. 25–36)
% Depends on macros from % jugeo-common.tex — shared preamble for all Judgment Geometry papers
% Include via: % jugeo-common.tex — shared preamble for all Judgment Geometry papers
% Include via: % jugeo-common.tex — shared preamble for all Judgment Geometry papers
% Include via: \input{jugeo-common}

% ── Packages ────────────────────────────────────────────────────────────
\usepackage{amsmath,amssymb,amsthm,mathtools}
\usepackage{tikz-cd}
\usepackage{listings}
\usepackage{xcolor}
\usepackage{xspace}
\usepackage{booktabs}
\usepackage{multirow}
\usepackage{enumitem}
\usepackage[expansion=false]{microtype}
\usepackage{hyperref}
\usepackage{cleveref}
% \usepackage{thmtools}  % disabled: not available in this TeX installation
\usepackage{float}
\usepackage{caption}
\usepackage{subcaption}
\usepackage{tabularx}
\usepackage{stmaryrd}       % \llbracket, \rrbracket
\usepackage{bbm}            % \mathbbm{1}

% ── Page geometry ───────────────────────────────────────────────────────
\usepackage[margin=1in]{geometry}

% ── Theorem environments ────────────────────────────────────────────────
\theoremstyle{plain}
\newtheorem{theorem}{Theorem}[section]
\newtheorem{lemma}[theorem]{Lemma}
\newtheorem{proposition}[theorem]{Proposition}
\newtheorem{corollary}[theorem]{Corollary}

\theoremstyle{definition}
\newtheorem{definition}[theorem]{Definition}
\newtheorem{example}[theorem]{Example}
\newtheorem{construction}[theorem]{Construction}

\theoremstyle{remark}
\newtheorem{remark}[theorem]{Remark}
\newtheorem{notation}[theorem]{Notation}

% ── Python listings style ──────────────────────────────────────────────
\definecolor{jugeoBlue}{HTML}{2563EB}
\definecolor{jugeoGreen}{HTML}{059669}
\definecolor{jugeoRed}{HTML}{DC2626}
\definecolor{jugeoPurple}{HTML}{7C3AED}
\definecolor{jugeoGray}{HTML}{6B7280}
\definecolor{jugeoBg}{HTML}{F8FAFC}

\lstdefinestyle{jugeo-python}{
  language=Python,
  basicstyle=\ttfamily\small,
  keywordstyle=\color{jugeoBlue}\bfseries,
  stringstyle=\color{jugeoGreen},
  commentstyle=\color{jugeoGray}\itshape,
  numberstyle=\tiny\color{jugeoGray},
  backgroundcolor=\color{jugeoBg},
  numbers=left,
  numbersep=6pt,
  frame=single,
  framerule=0.4pt,
  rulecolor=\color{jugeoGray!40},
  breaklines=true,
  breakatwhitespace=true,
  showstringspaces=false,
  tabsize=4,
  xleftmargin=1.5em,
  framexleftmargin=1.5em,
  morekeywords={dataclass,frozen,field,Enum,IntEnum,auto,Any,
                Mapping,Sequence,Optional,match,case},
  literate={→}{{$\to$}}1 {←}{{$\leftarrow$}}1
           {≤}{{$\leq$}}1 {≥}{{$\geq$}}1 {≠}{{$\neq$}}1
           {∈}{{$\in$}}1 {∀}{{$\forall$}}1 {∃}{{$\exists$}}1
           {⊕}{{$\oplus$}}1 {⊖}{{$\ominus$}}1,
  escapeinside={(*@}{@*)},
}
\lstset{style=jugeo-python}

\lstdefinestyle{jugeo-output}{
  basicstyle=\ttfamily\small,
  backgroundcolor=\color{jugeoBg},
  frame=single,
  framerule=0.4pt,
  rulecolor=\color{jugeoGray!40},
  breaklines=true,
  showstringspaces=false,
  xleftmargin=1.5em,
  framexleftmargin=0.5em,
  numbers=none,
}

% ── JuGeo-specific macros ──────────────────────────────────────────────

% Judgment 8-tuple
\newcommand{\J}{\mathcal{J}}
\newcommand{\Jud}[8]{(#1,\,#2,\,#3,\,#4,\,#5,\,#6,\,#7,\,#8)}
\newcommand{\JudShort}[1]{\J_{#1}}

% Components
\newcommand{\coord}{\mathsf{c}}            % coordinate
\newcommand{\prop}{\varphi}                 % proposition
\newcommand{\carrier}{\mathcal{A}}          % carrier type
\newcommand{\evid}{\mathcal{E}}             % evidence bundle
\newcommand{\oblig}{\mathcal{O}}            % obligations
\newcommand{\obstr}{\mathcal{B}}            % obstructions
\newcommand{\trust}{\mathcal{T}}            % trust annotation
\newcommand{\prov}{\Pi}                     % provenance

% Trust algebra
\newcommand{\Talg}{\mathfrak{T}}            % trust ordered algebra
\newcommand{\Eadm}{\mathcal{E}_{\mathrm{adm}}}
\newcommand{\tleq}{\preceq}                 % trust ordering
\newcommand{\tjoin}{\oplus}                 % conservative join
\newcommand{\tatten}{\ominus}               % attenuation
\newcommand{\tpromote}{\uparrow_\pi}        % promotion
\newcommand{\tdemote}{\downarrow_\chi}       % demotion

% Trust tiers
\newcommand{\Tcontra}{\ensuremath{\mathsf{CONTRADICTED}}}
\newcommand{\Tunver}{\ensuremath{\mathsf{UNVERIFIED}}}
\newcommand{\Tcopilot}{\ensuremath{\mathsf{COPILOT}}}
\newcommand{\Toracle}{\ensuremath{\mathsf{ORACLE}}}
\newcommand{\Truntime}{\ensuremath{\mathsf{RUNTIME}}}
\newcommand{\Tsolver}{\ensuremath{\mathsf{SOLVER}}}
\newcommand{\Tproof}{\ensuremath{\mathsf{PROOF}}}

% Site and geometry
\newcommand{\Site}{\mathcal{S}}             % semantic site
\newcommand{\Cov}{\mathrm{Cov}}            % covering family
\newcommand{\Ob}{\mathrm{Ob}}              % objects
\newcommand{\Mor}{\mathrm{Mor}}            % morphisms
\newcommand{\res}{\mathrm{res}}            % restriction
\newcommand{\Cech}{\check{C}}              % Čech complex
\newcommand{\Hcoh}[1]{H^{#1}}             % cohomology group

% Descent
\newcommand{\descent}{\mathrm{desc}}
\newcommand{\glue}{\mathrm{glue}}
\newcommand{\overlap}[2]{#1 \cap #2}

% Semantic moves
\newcommand{\VERIFY}{\textsc{Verify}}
\newcommand{\CONSTRUCT}{\textsc{Construct}}
\newcommand{\NORMALIZE}{\textsc{Normalize}}
\newcommand{\GLUE}{\textsc{Glue}}
\newcommand{\RESTRICT}{\textsc{Restrict}}
\newcommand{\TRANSPORT}{\textsc{Transport}}
\newcommand{\REPAIR}{\textsc{Repair}}
\newcommand{\PROMOTE}{\textsc{Promote}}
\newcommand{\CHALLENGE}{\textsc{Challenge}}
\newcommand{\ESCALATE}{\textsc{Escalate}}

% Fragments
\newcommand{\QFLIA}{\texttt{QF\_LIA}}
\newcommand{\QFLRA}{\texttt{QF\_LRA}}
\newcommand{\QFBV}{\texttt{QF\_BV}}
\newcommand{\QFUF}{\texttt{QF\_UF}}

% Misc
\newcommand{\Python}{\textsc{Python}}
\newcommand{\JuGeo}{\textsc{JuGeo}}
\newcommand{\LEAN}{\textsc{Lean}\,4}
\newcommand{\Fstar}{F$^{\star}$}
\newcommand{\Coq}{\textsc{Coq}}
\newcommand{\Dafny}{\textsc{Dafny}}

% Common shortcuts
\newcommand{\ie}{i.e.\@\xspace}
\newcommand{\eg}{e.g.\@\xspace}
\newcommand{\cf}{cf.\@\xspace}
\newcommand{\wrt}{w.r.t.\@\xspace}

% Evaluation shortcuts
\newcommand{\numcases}{300}
\newcommand{\numspec}{100}
\newcommand{\numeq}{100}
\newcommand{\numbug}{100}
\newcommand{\medtime}{6.5\,ms}
\newcommand{\totaltime}{1.949\,s}


% ── Packages ────────────────────────────────────────────────────────────
\usepackage{amsmath,amssymb,amsthm,mathtools}
\usepackage{tikz-cd}
\usepackage{listings}
\usepackage{xcolor}
\usepackage{xspace}
\usepackage{booktabs}
\usepackage{multirow}
\usepackage{enumitem}
\usepackage[expansion=false]{microtype}
\usepackage{hyperref}
\usepackage{cleveref}
% \usepackage{thmtools}  % disabled: not available in this TeX installation
\usepackage{float}
\usepackage{caption}
\usepackage{subcaption}
\usepackage{tabularx}
\usepackage{stmaryrd}       % \llbracket, \rrbracket
\usepackage{bbm}            % \mathbbm{1}

% ── Page geometry ───────────────────────────────────────────────────────
\usepackage[margin=1in]{geometry}

% ── Theorem environments ────────────────────────────────────────────────
\theoremstyle{plain}
\newtheorem{theorem}{Theorem}[section]
\newtheorem{lemma}[theorem]{Lemma}
\newtheorem{proposition}[theorem]{Proposition}
\newtheorem{corollary}[theorem]{Corollary}

\theoremstyle{definition}
\newtheorem{definition}[theorem]{Definition}
\newtheorem{example}[theorem]{Example}
\newtheorem{construction}[theorem]{Construction}

\theoremstyle{remark}
\newtheorem{remark}[theorem]{Remark}
\newtheorem{notation}[theorem]{Notation}

% ── Python listings style ──────────────────────────────────────────────
\definecolor{jugeoBlue}{HTML}{2563EB}
\definecolor{jugeoGreen}{HTML}{059669}
\definecolor{jugeoRed}{HTML}{DC2626}
\definecolor{jugeoPurple}{HTML}{7C3AED}
\definecolor{jugeoGray}{HTML}{6B7280}
\definecolor{jugeoBg}{HTML}{F8FAFC}

\lstdefinestyle{jugeo-python}{
  language=Python,
  basicstyle=\ttfamily\small,
  keywordstyle=\color{jugeoBlue}\bfseries,
  stringstyle=\color{jugeoGreen},
  commentstyle=\color{jugeoGray}\itshape,
  numberstyle=\tiny\color{jugeoGray},
  backgroundcolor=\color{jugeoBg},
  numbers=left,
  numbersep=6pt,
  frame=single,
  framerule=0.4pt,
  rulecolor=\color{jugeoGray!40},
  breaklines=true,
  breakatwhitespace=true,
  showstringspaces=false,
  tabsize=4,
  xleftmargin=1.5em,
  framexleftmargin=1.5em,
  morekeywords={dataclass,frozen,field,Enum,IntEnum,auto,Any,
                Mapping,Sequence,Optional,match,case},
  literate={→}{{$\to$}}1 {←}{{$\leftarrow$}}1
           {≤}{{$\leq$}}1 {≥}{{$\geq$}}1 {≠}{{$\neq$}}1
           {∈}{{$\in$}}1 {∀}{{$\forall$}}1 {∃}{{$\exists$}}1
           {⊕}{{$\oplus$}}1 {⊖}{{$\ominus$}}1,
  escapeinside={(*@}{@*)},
}
\lstset{style=jugeo-python}

\lstdefinestyle{jugeo-output}{
  basicstyle=\ttfamily\small,
  backgroundcolor=\color{jugeoBg},
  frame=single,
  framerule=0.4pt,
  rulecolor=\color{jugeoGray!40},
  breaklines=true,
  showstringspaces=false,
  xleftmargin=1.5em,
  framexleftmargin=0.5em,
  numbers=none,
}

% ── JuGeo-specific macros ──────────────────────────────────────────────

% Judgment 8-tuple
\newcommand{\J}{\mathcal{J}}
\newcommand{\Jud}[8]{(#1,\,#2,\,#3,\,#4,\,#5,\,#6,\,#7,\,#8)}
\newcommand{\JudShort}[1]{\J_{#1}}

% Components
\newcommand{\coord}{\mathsf{c}}            % coordinate
\newcommand{\prop}{\varphi}                 % proposition
\newcommand{\carrier}{\mathcal{A}}          % carrier type
\newcommand{\evid}{\mathcal{E}}             % evidence bundle
\newcommand{\oblig}{\mathcal{O}}            % obligations
\newcommand{\obstr}{\mathcal{B}}            % obstructions
\newcommand{\trust}{\mathcal{T}}            % trust annotation
\newcommand{\prov}{\Pi}                     % provenance

% Trust algebra
\newcommand{\Talg}{\mathfrak{T}}            % trust ordered algebra
\newcommand{\Eadm}{\mathcal{E}_{\mathrm{adm}}}
\newcommand{\tleq}{\preceq}                 % trust ordering
\newcommand{\tjoin}{\oplus}                 % conservative join
\newcommand{\tatten}{\ominus}               % attenuation
\newcommand{\tpromote}{\uparrow_\pi}        % promotion
\newcommand{\tdemote}{\downarrow_\chi}       % demotion

% Trust tiers
\newcommand{\Tcontra}{\ensuremath{\mathsf{CONTRADICTED}}}
\newcommand{\Tunver}{\ensuremath{\mathsf{UNVERIFIED}}}
\newcommand{\Tcopilot}{\ensuremath{\mathsf{COPILOT}}}
\newcommand{\Toracle}{\ensuremath{\mathsf{ORACLE}}}
\newcommand{\Truntime}{\ensuremath{\mathsf{RUNTIME}}}
\newcommand{\Tsolver}{\ensuremath{\mathsf{SOLVER}}}
\newcommand{\Tproof}{\ensuremath{\mathsf{PROOF}}}

% Site and geometry
\newcommand{\Site}{\mathcal{S}}             % semantic site
\newcommand{\Cov}{\mathrm{Cov}}            % covering family
\newcommand{\Ob}{\mathrm{Ob}}              % objects
\newcommand{\Mor}{\mathrm{Mor}}            % morphisms
\newcommand{\res}{\mathrm{res}}            % restriction
\newcommand{\Cech}{\check{C}}              % Čech complex
\newcommand{\Hcoh}[1]{H^{#1}}             % cohomology group

% Descent
\newcommand{\descent}{\mathrm{desc}}
\newcommand{\glue}{\mathrm{glue}}
\newcommand{\overlap}[2]{#1 \cap #2}

% Semantic moves
\newcommand{\VERIFY}{\textsc{Verify}}
\newcommand{\CONSTRUCT}{\textsc{Construct}}
\newcommand{\NORMALIZE}{\textsc{Normalize}}
\newcommand{\GLUE}{\textsc{Glue}}
\newcommand{\RESTRICT}{\textsc{Restrict}}
\newcommand{\TRANSPORT}{\textsc{Transport}}
\newcommand{\REPAIR}{\textsc{Repair}}
\newcommand{\PROMOTE}{\textsc{Promote}}
\newcommand{\CHALLENGE}{\textsc{Challenge}}
\newcommand{\ESCALATE}{\textsc{Escalate}}

% Fragments
\newcommand{\QFLIA}{\texttt{QF\_LIA}}
\newcommand{\QFLRA}{\texttt{QF\_LRA}}
\newcommand{\QFBV}{\texttt{QF\_BV}}
\newcommand{\QFUF}{\texttt{QF\_UF}}

% Misc
\newcommand{\Python}{\textsc{Python}}
\newcommand{\JuGeo}{\textsc{JuGeo}}
\newcommand{\LEAN}{\textsc{Lean}\,4}
\newcommand{\Fstar}{F$^{\star}$}
\newcommand{\Coq}{\textsc{Coq}}
\newcommand{\Dafny}{\textsc{Dafny}}

% Common shortcuts
\newcommand{\ie}{i.e.\@\xspace}
\newcommand{\eg}{e.g.\@\xspace}
\newcommand{\cf}{cf.\@\xspace}
\newcommand{\wrt}{w.r.t.\@\xspace}

% Evaluation shortcuts
\newcommand{\numcases}{300}
\newcommand{\numspec}{100}
\newcommand{\numeq}{100}
\newcommand{\numbug}{100}
\newcommand{\medtime}{6.5\,ms}
\newcommand{\totaltime}{1.949\,s}


% ── Packages ────────────────────────────────────────────────────────────
\usepackage{amsmath,amssymb,amsthm,mathtools}
\usepackage{tikz-cd}
\usepackage{listings}
\usepackage{xcolor}
\usepackage{xspace}
\usepackage{booktabs}
\usepackage{multirow}
\usepackage{enumitem}
\usepackage[expansion=false]{microtype}
\usepackage{hyperref}
\usepackage{cleveref}
% \usepackage{thmtools}  % disabled: not available in this TeX installation
\usepackage{float}
\usepackage{caption}
\usepackage{subcaption}
\usepackage{tabularx}
\usepackage{stmaryrd}       % \llbracket, \rrbracket
\usepackage{bbm}            % \mathbbm{1}

% ── Page geometry ───────────────────────────────────────────────────────
\usepackage[margin=1in]{geometry}

% ── Theorem environments ────────────────────────────────────────────────
\theoremstyle{plain}
\newtheorem{theorem}{Theorem}[section]
\newtheorem{lemma}[theorem]{Lemma}
\newtheorem{proposition}[theorem]{Proposition}
\newtheorem{corollary}[theorem]{Corollary}

\theoremstyle{definition}
\newtheorem{definition}[theorem]{Definition}
\newtheorem{example}[theorem]{Example}
\newtheorem{construction}[theorem]{Construction}

\theoremstyle{remark}
\newtheorem{remark}[theorem]{Remark}
\newtheorem{notation}[theorem]{Notation}

% ── Python listings style ──────────────────────────────────────────────
\definecolor{jugeoBlue}{HTML}{2563EB}
\definecolor{jugeoGreen}{HTML}{059669}
\definecolor{jugeoRed}{HTML}{DC2626}
\definecolor{jugeoPurple}{HTML}{7C3AED}
\definecolor{jugeoGray}{HTML}{6B7280}
\definecolor{jugeoBg}{HTML}{F8FAFC}

\lstdefinestyle{jugeo-python}{
  language=Python,
  basicstyle=\ttfamily\small,
  keywordstyle=\color{jugeoBlue}\bfseries,
  stringstyle=\color{jugeoGreen},
  commentstyle=\color{jugeoGray}\itshape,
  numberstyle=\tiny\color{jugeoGray},
  backgroundcolor=\color{jugeoBg},
  numbers=left,
  numbersep=6pt,
  frame=single,
  framerule=0.4pt,
  rulecolor=\color{jugeoGray!40},
  breaklines=true,
  breakatwhitespace=true,
  showstringspaces=false,
  tabsize=4,
  xleftmargin=1.5em,
  framexleftmargin=1.5em,
  morekeywords={dataclass,frozen,field,Enum,IntEnum,auto,Any,
                Mapping,Sequence,Optional,match,case},
  literate={→}{{$\to$}}1 {←}{{$\leftarrow$}}1
           {≤}{{$\leq$}}1 {≥}{{$\geq$}}1 {≠}{{$\neq$}}1
           {∈}{{$\in$}}1 {∀}{{$\forall$}}1 {∃}{{$\exists$}}1
           {⊕}{{$\oplus$}}1 {⊖}{{$\ominus$}}1,
  escapeinside={(*@}{@*)},
}
\lstset{style=jugeo-python}

\lstdefinestyle{jugeo-output}{
  basicstyle=\ttfamily\small,
  backgroundcolor=\color{jugeoBg},
  frame=single,
  framerule=0.4pt,
  rulecolor=\color{jugeoGray!40},
  breaklines=true,
  showstringspaces=false,
  xleftmargin=1.5em,
  framexleftmargin=0.5em,
  numbers=none,
}

% ── JuGeo-specific macros ──────────────────────────────────────────────

% Judgment 8-tuple
\newcommand{\J}{\mathcal{J}}
\newcommand{\Jud}[8]{(#1,\,#2,\,#3,\,#4,\,#5,\,#6,\,#7,\,#8)}
\newcommand{\JudShort}[1]{\J_{#1}}

% Components
\newcommand{\coord}{\mathsf{c}}            % coordinate
\newcommand{\prop}{\varphi}                 % proposition
\newcommand{\carrier}{\mathcal{A}}          % carrier type
\newcommand{\evid}{\mathcal{E}}             % evidence bundle
\newcommand{\oblig}{\mathcal{O}}            % obligations
\newcommand{\obstr}{\mathcal{B}}            % obstructions
\newcommand{\trust}{\mathcal{T}}            % trust annotation
\newcommand{\prov}{\Pi}                     % provenance

% Trust algebra
\newcommand{\Talg}{\mathfrak{T}}            % trust ordered algebra
\newcommand{\Eadm}{\mathcal{E}_{\mathrm{adm}}}
\newcommand{\tleq}{\preceq}                 % trust ordering
\newcommand{\tjoin}{\oplus}                 % conservative join
\newcommand{\tatten}{\ominus}               % attenuation
\newcommand{\tpromote}{\uparrow_\pi}        % promotion
\newcommand{\tdemote}{\downarrow_\chi}       % demotion

% Trust tiers
\newcommand{\Tcontra}{\ensuremath{\mathsf{CONTRADICTED}}}
\newcommand{\Tunver}{\ensuremath{\mathsf{UNVERIFIED}}}
\newcommand{\Tcopilot}{\ensuremath{\mathsf{COPILOT}}}
\newcommand{\Toracle}{\ensuremath{\mathsf{ORACLE}}}
\newcommand{\Truntime}{\ensuremath{\mathsf{RUNTIME}}}
\newcommand{\Tsolver}{\ensuremath{\mathsf{SOLVER}}}
\newcommand{\Tproof}{\ensuremath{\mathsf{PROOF}}}

% Site and geometry
\newcommand{\Site}{\mathcal{S}}             % semantic site
\newcommand{\Cov}{\mathrm{Cov}}            % covering family
\newcommand{\Ob}{\mathrm{Ob}}              % objects
\newcommand{\Mor}{\mathrm{Mor}}            % morphisms
\newcommand{\res}{\mathrm{res}}            % restriction
\newcommand{\Cech}{\check{C}}              % Čech complex
\newcommand{\Hcoh}[1]{H^{#1}}             % cohomology group

% Descent
\newcommand{\descent}{\mathrm{desc}}
\newcommand{\glue}{\mathrm{glue}}
\newcommand{\overlap}[2]{#1 \cap #2}

% Semantic moves
\newcommand{\VERIFY}{\textsc{Verify}}
\newcommand{\CONSTRUCT}{\textsc{Construct}}
\newcommand{\NORMALIZE}{\textsc{Normalize}}
\newcommand{\GLUE}{\textsc{Glue}}
\newcommand{\RESTRICT}{\textsc{Restrict}}
\newcommand{\TRANSPORT}{\textsc{Transport}}
\newcommand{\REPAIR}{\textsc{Repair}}
\newcommand{\PROMOTE}{\textsc{Promote}}
\newcommand{\CHALLENGE}{\textsc{Challenge}}
\newcommand{\ESCALATE}{\textsc{Escalate}}

% Fragments
\newcommand{\QFLIA}{\texttt{QF\_LIA}}
\newcommand{\QFLRA}{\texttt{QF\_LRA}}
\newcommand{\QFBV}{\texttt{QF\_BV}}
\newcommand{\QFUF}{\texttt{QF\_UF}}

% Misc
\newcommand{\Python}{\textsc{Python}}
\newcommand{\JuGeo}{\textsc{JuGeo}}
\newcommand{\LEAN}{\textsc{Lean}\,4}
\newcommand{\Fstar}{F$^{\star}$}
\newcommand{\Coq}{\textsc{Coq}}
\newcommand{\Dafny}{\textsc{Dafny}}

% Common shortcuts
\newcommand{\ie}{i.e.\@\xspace}
\newcommand{\eg}{e.g.\@\xspace}
\newcommand{\cf}{cf.\@\xspace}
\newcommand{\wrt}{w.r.t.\@\xspace}

% Evaluation shortcuts
\newcommand{\numcases}{300}
\newcommand{\numspec}{100}
\newcommand{\numeq}{100}
\newcommand{\numbug}{100}
\newcommand{\medtime}{6.5\,ms}
\newcommand{\totaltime}{1.949\,s}


%% ═══════════════════════════════════════════════════════════════════════
\section{The Semantic Moves Calculus}\label{sec:moves}
%% ═══════════════════════════════════════════════════════════════════════

The preceding sections established a static framework: semantic sites
furnish a category of program coordinates, judgments live in presheaves
over those sites, and obstructions in the \v{C}ech cohomology
classify why local verdicts fail to glue.  We now give the framework
a \emph{dynamics} by defining \textbf{semantic moves}---typed,
composable operations that transform proof states while preserving
or increasing trust.

The key insight is that verification is not a single monolithic
decision but a \emph{sequence of categorical transformations}: restrict
a judgment to a local coordinate, verify it there, glue the local
verdicts back, and---if an obstruction is found---repair the program
so the obstruction vanishes.  Each step is a morphism in a category of
proof states, and compositionality of morphisms gives us compositionality
of the verification pipeline.

\subsection{Proof States and the Move Category}\label{sec:move-category}

\begin{definition}[Proof State]\label{def:seminal-proof-state}
A \emph{proof state} is a tuple
\[
  \mathcal{P} \;=\; \bigl(\Site_P,\; \J_\bullet,\;
    \evid_\bullet,\; \oblig_\bullet,\; \obstr_\bullet,\;
    \trust_\bullet,\; \prov_\bullet\bigr)
\]
where $\Site_P = (\mathcal{C}_P, J_P)$ is the semantic site
(\Cref{sec:sites}); $\J_\bullet$ is a presheaf of judgments;
$\evid_\bullet$, $\oblig_\bullet$, $\obstr_\bullet$ are presheaves of
evidence bundles, obligation sets, and obstruction cocycles;
$\trust_\bullet : \Ob(\mathcal{C}_P) \to \Talg$ assigns trust levels;
and $\prov_\bullet$ records provenance.  A proof state is
\emph{well-formed} if:
\begin{enumerate}[label=\textnormal{(WF\arabic*)},nosep]
  \item every judgment satisfies
        $\trust_\bullet(c) \tleq \trust(\J_c)$;
  \item for every morphism $f : c \to c'$,
        $\res_f(\J_{c'}) = \J_c$;
  \item $\obstr_\bullet$ is a \v{C}ech 1-cocycle.
\end{enumerate}
\end{definition}

\begin{definition}[Move Category]\label{def:move-cat}
The \emph{move category} $\mathbf{M}$ has well-formed proof states as
objects and semantic moves $m : \mathcal{P} \to \mathcal{P}'$ as
morphisms, where $m$ sends well-formed states to well-formed states.
Identity is the no-op move; composition is sequential application.
\end{definition}

\subsection{Move Definitions}\label{sec:move-defs}

Each semantic move is specified by a precondition, a transformation
rule, and a trust effect.  We define ten canonical moves.

\begin{definition}[Canonical Semantic Moves]\label{def:canonical-moves}
Let $\mathcal{P}$ be a well-formed proof state.
\end{definition}

\paragraph{\VERIFY.}
Check a proposition $\prop$ at coordinate $c$, producing an evidence
item $e$.
\begin{itemize}[nosep]
  \item \emph{Precondition:} $c \in \Ob(\mathcal{C}_P)$ and $\prop$
        is a well-formed proposition at $c$.
  \item \emph{Transformation:}
        $\evid'_c = \evid_c \cup \{e\}$, where $e$ is obtained from
        the appropriate evidence channel (SMT solver, runtime, oracle).
  \item \emph{Trust effect:}
        $\trust'(c) = \trust(c) \tjoin \trust(e)$.
\end{itemize}
The $\VERIFY$ move is the fundamental evidence-producing operation.
The evidence item $e$ carries a trust level determined by its
source channel: $\Tsolver$ for an SMT proof, $\Truntime$ for a
test witness, $\Tcopilot$ for an LLM suggestion.

\paragraph{\CONSTRUCT.}
Synthesize a witness for an existential proposition
$\exists x.\, \prop(x)$ at coordinate $c$.
\begin{itemize}[nosep]
  \item \emph{Precondition:}
        $\prop_c = \exists x.\, \psi(x)$ for some $\psi$.
  \item \emph{Transformation:}
        Find a term $w$ such that $\psi(w)$ holds; add evidence
        $e_w = (\mathsf{WITNESS}, w, \trust_w)$ to $\evid_c$;
        replace the proposition with $\psi(w)$.
  \item \emph{Trust effect:} $\trust'(c) = \trust(c) \tjoin \trust(e_w)$.
\end{itemize}

\paragraph{\NORMALIZE.}
Reduce a judgment $\J_c$ to canonical form by applying algebraic
simplifications.
\begin{itemize}[nosep]
  \item \emph{Precondition:} $\J_c$ is well-formed.
  \item \emph{Transformation:}
        Apply the five judgment operations (restriction, transport,
        composition, comparison, join) to produce an equivalent judgment
        $\J'_c$ in normal form: discharged obligations removed,
        evidence deduplicated, carrier simplified.
  \item \emph{Trust effect:} unchanged ($\trust'(c) = \trust(c)$).
\end{itemize}

\paragraph{\GLUE.}
Combine compatible local judgments into a global judgment---the
descent step.
\begin{itemize}[nosep]
  \item \emph{Precondition:}
        A covering family $\{c_i \to c\}_{i \in I}$ and local
        sections $\{s_i = \J_{c_i}\}_{i \in I}$ satisfying the
        cocycle condition:
        $\res_{c_i \cap c_j \hookrightarrow c_i}(s_i) =
         \res_{c_i \cap c_j \hookrightarrow c_j}(s_j)$ for all
        $i,j \in I$.
  \item \emph{Transformation:}
        Produce a global section $s \in \J(c)$ with
        $\res_{c_i \hookrightarrow c}(s) = s_i$ for all $i$.
  \item \emph{Trust effect:}
        $\trust'(c) = \bigoplus_{i \in I} \trust(c_i)$
        (conservative join over all patches).
\end{itemize}
The $\GLUE$ move is the cornerstone of the pipeline.  When it
succeeds, the program has been verified; when it fails, the failure
is recorded as an obstruction (\cf \REPAIR).

\paragraph{\RESTRICT.}
Pull back a judgment along a morphism $f : U \to V$.
\begin{itemize}[nosep]
  \item \emph{Precondition:}
        $f \in \Mor(\mathcal{C}_P)$ and $\J_V$ is defined.
  \item \emph{Transformation:}
        $\J'_U = \res_f(\J_V)$, with evidence and obligations
        pulled back componentwise.
  \item \emph{Trust effect:}
        $\trust'(U) = \trust(V) \tatten \delta(f)$, where
        $\delta(f)$ is the attenuation factor determined by the
        morphism kind (restriction morphisms attenuate less than
        transport morphisms).
\end{itemize}

\paragraph{\TRANSPORT.}
Move a judgment along an equivalence $f : c \xrightarrow{\sim} c'$.
\begin{itemize}[nosep]
  \item \emph{Precondition:}
        $f$ is invertible in $\mathcal{C}_P$.
  \item \emph{Transformation:}
        $\J'_{c'} = f_*(\J_c)$, the pushforward along $f$.
  \item \emph{Trust effect:}
        $\trust'(c') = \trust(c)$ (no attenuation for equivalences).
\end{itemize}

\paragraph{\REPAIR.}
Given an obstruction $[\alpha] \in \Hcoh{1}(\Cech(\mathcal{U}), \J)$,
modify the program to kill $[\alpha]$.
\begin{itemize}[nosep]
  \item \emph{Precondition:}
        $\Hcoh{1} \neq 0$ and $[\alpha]$ is classified with repair
        hints.
  \item \emph{Transformation:}
        Localize $[\alpha]$ to its support coordinate $c_\alpha$;
        apply a code modification $\sigma$ at $c_\alpha$ such that
        $[\alpha]$ becomes a coboundary in the modified program;
        recompute the affected local sections.
  \item \emph{Trust effect:}
        $\trust'(c_\alpha) = \Tcopilot$ initially (repair is a
        suggestion until re-verified).
\end{itemize}

\paragraph{\PROMOTE.}
Upgrade a trust level given new evidence.
\begin{itemize}[nosep]
  \item \emph{Precondition:}
        New evidence $e$ at $c$ with $\trust(e) \succ \trust(c)$.
  \item \emph{Transformation:}
        $\trust'(c) = \tpromote(\trust(c), \trust(e))$.
  \item \emph{Trust effect:} strictly increasing.
\end{itemize}

\paragraph{\CHALLENGE.}
Attempt to refute a judgment (the adversarial move).
\begin{itemize}[nosep]
  \item \emph{Precondition:} $\J_c$ is defined.
  \item \emph{Transformation:}
        Search for a counterexample to $\prop_c$.  If found, record
        an obstruction and set $\trust'(c) = \Tcontra$.  If not
        found after bounded search, the judgment stands.
  \item \emph{Trust effect:}
        $\trust'(c) = \Tcontra$ on success (demotion);
        unchanged on failure to refute.
\end{itemize}

\paragraph{\ESCALATE.}
Raise a judgment to a higher trust tier by invoking a stronger oracle.
\begin{itemize}[nosep]
  \item \emph{Precondition:}
        $\trust(c) \prec \Tproof$ and a stronger evidence channel
        is available.
  \item \emph{Transformation:}
        Invoke the stronger channel (e.g.\ upgrade from runtime
        witness to SMT proof); add the resulting evidence to
        $\evid_c$.
  \item \emph{Trust effect:}
        $\trust'(c) = \trust(c) \tjoin \trust(e_{\mathrm{new}})$
        (non-decreasing).
\end{itemize}

\begin{remark}[Classical vs.\ Geometric Moves]\label{rem:class-geo}
The moves $\VERIFY$, $\CONSTRUCT$, $\NORMALIZE$, $\PROMOTE$,
$\CHALLENGE$, and $\ESCALATE$ are \emph{classical}: they have
analogues in tactic-based proof assistants (analogous to
\texttt{apply}, \texttt{exact}, \texttt{simp}, and
\texttt{omega} in \LEAN).  The moves $\GLUE$, $\RESTRICT$,
$\TRANSPORT$, and $\REPAIR$ are \emph{geometric}: they act on
the site structure (covers, overlaps, cohomology) and have no
analogue in \LEAN, \Coq, or \Fstar.  This geometric fraction is
the source of \JuGeo's additional expressive power.
\end{remark}

\subsection{Move Composition}\label{sec:move-composition}

Moves compose in two ways: sequentially and in parallel.

\begin{definition}[Sequential Composition]\label{def:seq-compose}
Given $m_1 : \mathcal{P}_0 \to \mathcal{P}_1$ and
$m_2 : \mathcal{P}_1 \to \mathcal{P}_2$, define
$m_2 \circ m_1 : \mathcal{P}_0 \to \mathcal{P}_2$ by
$(m_2 \circ m_1)(\mathcal{P}_0) = m_2(m_1(\mathcal{P}_0))$.
\end{definition}

\begin{definition}[Parallel Composition]\label{def:par-compose}
When $m_1$ and $m_2$ act on disjoint coordinates, define
$m_1 \otimes m_2$ by applying both simultaneously.  The trust of
the composite is the conservative join:
$\trust(m_1 \otimes m_2) = \trust(m_1) \tjoin \trust(m_2)$.
\end{definition}

Canonical composite moves realize the major patterns of the pipeline:

\begin{enumerate}[label=(\roman*),nosep]
  \item \textbf{Local verification:}
    $\VERIFY \circ \RESTRICT$ restricts a judgment to a local
    coordinate and verifies it there.

  \item \textbf{Full descent pipeline:}
    $\GLUE \circ \bigl(\VERIFY_{c_1} \otimes \cdots \otimes
     \VERIFY_{c_n}\bigr) \circ \bigl(\RESTRICT_{c_1} \otimes
     \cdots \otimes \RESTRICT_{c_n}\bigr)$
    restricts, verifies each patch in parallel, and glues.

  \item \textbf{Adversarial repair:}
    $\REPAIR \circ \CHALLENGE$ first attempts to refute and then
    repairs any obstructions found.

  \item \textbf{Trust escalation:}
    $\ESCALATE \circ \VERIFY$ verifies at a low tier and then
    upgrades the evidence.

  \item \textbf{Iterative refinement:}
    $(\GLUE \circ \VERIFY^{\otimes n} \circ \RESTRICT^{\otimes n})
    \circ \REPAIR$ repairs, re-restricts, re-verifies, and re-glues.
\end{enumerate}

\subsection{Soundness Theorem}\label{sec:move-soundness}

\begin{theorem}[Move Soundness]\label{thm:move-soundness}
Every semantic move preserves well-formedness and is monotone
in trust.  Formally: if move $m$ transforms proof state
$\mathcal{P}$ to $\mathcal{P}'$, then
\begin{enumerate}[label=\textnormal{(\alph*)},nosep]
  \item $\mathcal{P}'$ is well-formed (conditions
        \textnormal{WF1--WF3} hold);
  \item either $\trust(\mathcal{P}') \succeq \trust(\mathcal{P})$,
        or $m$ is a $\CHALLENGE$ that discovered a genuine refutation.
\end{enumerate}
\end{theorem}

\begin{proof}
By case analysis on the move kind.

\emph{Evidence-producing moves} ($\VERIFY$, $\CONSTRUCT$,
$\PROMOTE$, $\ESCALATE$): these add evidence to $\evid_c$ and
apply the conservative join $\tjoin$ to the trust level.  Since
$\tjoin$ is monotone and $\evid$ is a presheaf (restriction
commutes with union), WF1 and WF2 are preserved.  WF3 is
unaffected because the obstruction presheaf is unchanged.

\emph{Normalization} ($\NORMALIZE$): rewrites the judgment at
a single fiber.  The rewrite preserves the judgment's denotation
(it is an algebraic identity), so WF2 is maintained.  Trust is
unchanged.

\emph{Restriction and transport} ($\RESTRICT$, $\TRANSPORT$):
by functoriality of the presheaf $\J_\bullet$, restriction maps
compose with existing restrictions, preserving WF2.  For
$\TRANSPORT$, invertibility of $f$ ensures that the pushforward
$f_*$ is an isomorphism on fibers, so no information is lost and
trust is preserved exactly.  For $\RESTRICT$, the attenuation
factor $\delta(f) \geq 0$ ensures $\trust' \tleq \trust$, but
since $\tleq$ is the trust ordering, the trust of the restricted
judgment is bounded by the original---WF1 is preserved.

\emph{Gluing} ($\GLUE$): the precondition requires the cocycle
condition.  When the cocycle condition holds, the descent theorem
(\Cref{sec:descent}) guarantees a unique global section satisfying
WF2.  The trust of the global section is the conservative join
$\bigoplus_i \trust(c_i)$, which is $\tleq$ every $\trust(c_i)$
(since $\tjoin$ takes the minimum), preserving WF1.  The
obstruction presheaf is set to the trivial cocycle ($\Hcoh{1} = 0$),
preserving WF3.

\emph{Repair} ($\REPAIR$): modifies the site by adjusting code at
a support coordinate.  The new covering family satisfies the
Grothendieck axioms by pullback stability (Axiom T2 from
\Cref{sec:sites}).  Local sections are recomputed, so WF2 holds
for the new site.  Trust is reset to $\Tcopilot$ at the repaired
coordinate, which is admissible, so WF1 holds.  The obstruction
class $[\alpha]$ is removed from $\obstr_\bullet$ (it has become
a coboundary), maintaining WF3.

\emph{Challenge} ($\CHALLENGE$): if a counterexample is found,
$\trust'(c) = \Tcontra$ and an obstruction is added.  The
obstruction is a valid 1-cocycle (it witnesses the overlap
failure), so WF3 holds.  This is the unique move that can
\emph{decrease} trust, and it does so only when a genuine
refutation is produced.  If no counterexample is found, the proof
state is unchanged.
\end{proof}

\begin{corollary}[No Trust Laundering]\label{cor:no-laundering}
No finite composition of semantic moves can promote a judgment from
$\Tunver$ to $\Tsolver$ without producing solver-level evidence.
\end{corollary}

\begin{proof}
By induction on the length of the move trace.  The only moves that
increase trust are $\VERIFY$, $\CONSTRUCT$, $\PROMOTE$, and
$\ESCALATE$, all of which require evidence at or above the target
trust level.  The conservative join $\tjoin$ cannot exceed the trust
of the strongest evidence item.  Therefore, to reach $\Tsolver$,
some evidence item $e$ with $\trust(e) \geq \Tsolver$ must have
been produced, which requires an SMT solver proof.
\end{proof}


%% ═══════════════════════════════════════════════════════════════════════
\section{SMT Dispatch and Fragment Routing}\label{sec:dispatch}
%% ═══════════════════════════════════════════════════════════════════════

The $\VERIFY$ move relies on an evidence channel to produce a verdict.
For propositions in decidable first-order fragments, the evidence
channel is an SMT solver.  This section describes how \JuGeo{}
classifies each proposition into the most specific SMT-LIB fragment and
dispatches it to Z3 with optimal theory configuration, keeping solving
time minimal while maintaining encoding soundness.

\subsection{Fragment Classification}\label{sec:fragment-class}

\JuGeo{} maintains a taxonomy of 14 SMT-LIB fragments organized as a
bounded lattice $(\mathcal{F}, \sqsubseteq)$ by theory inclusion.
The four principal quantifier-free fragments are:

\begin{center}
\begin{tabular}{@{}llll@{}}
\toprule
\textbf{Fragment} & \textbf{Theory} & \textbf{Decidable?}
  & \textbf{Complexity} \\
\midrule
\QFLIA  & Linear integer arithmetic   & Yes & NP-complete \\
\QFLRA  & Linear real arithmetic      & Yes & Polynomial \\
\QFBV   & Fixed-width bitvectors      & Yes & NP-complete \\
\QFUF   & Uninterpreted functions     & Yes & $O(n \log n)$ \\
\midrule
\texttt{QF\_UFLIA} & UF + LIA        & Yes & NP-complete \\
\texttt{QF\_AUFLIA} & Arrays + UF + LIA & Yes & NP-complete \\
\bottomrule
\end{tabular}
\end{center}

\noindent
The lattice has bottom $\bot = \QFUF$ (the most restrictive decidable
fragment) and top $\top = \texttt{UNKNOWN}$.  An arrow
$F_1 \sqsubseteq F_2$ means every formula in $F_1$ is also in $F_2$.

\begin{definition}[Fragment Membership]\label{def:frag-member}
A closed formula $\varphi$ \emph{belongs to} fragment $F$, written
$\varphi \in F$, if (a) every non-logical symbol in $\varphi$ belongs
to the signature of $F$, (b) the syntactic restrictions of $F$ are
satisfied (\eg quantifier-free, linear), and (c) no proper
sub-fragment $F' \sqsubset F$ satisfies (a) and (b).  Condition~(c)
ensures that membership returns the \emph{most specific} fragment.
\end{definition}

The classification algorithm extracts a \emph{syntactic signature}
from each proposition in a single pass over its AST:

\begin{lstlisting}[style=jugeo-python,
  caption={Fragment classification (simplified).},
  label={lst:classify}]
def classify_fragment(phi: SMTFormula) -> Fragment:
    sig = extract_signature(phi)  # O(|phi|) pass
    if sig.has_quantifiers:
        return Fragment.QUANTIFIED
    if sig.has_arrays and sig.has_bitvectors:
        return Fragment.QF_ABV
    if sig.has_arrays and sig.has_integers:
        return Fragment.QF_AUFLIA
    if sig.has_uninterpreted and sig.has_integers:
        return Fragment.QF_UFLIA
    if sig.has_bitvectors:
        return Fragment.QF_BV
    if sig.has_integers and not sig.has_nonlinear:
        return Fragment.QF_LIA
    if sig.has_reals and not sig.has_nonlinear:
        return Fragment.QF_LRA
    if sig.has_uninterpreted:
        return Fragment.QF_UF
    return Fragment.UNKNOWN
\end{lstlisting}

\begin{lemma}[Classification Correctness]\label{lem:classify}
The function \texttt{classify\_fragment} runs in $O(|\varphi|)$ time
and returns the most specific fragment $F$ such that $\varphi \in F$.
\end{lemma}

\begin{proof}
The signature extraction performs a single traversal of the formula
AST, recording which sorts (integer, real, bitvector), function
symbols (arrays, uninterpreted), and syntactic features (quantifiers,
nonlinear multiplication) appear.  The cascade of conditionals tests
fragments in order from most general to most specific.
Condition~(c) of \Cref{def:frag-member} is satisfied because we
return the first (most specific) match.
\end{proof}

\subsection{The Dispatch Algorithm}\label{sec:dispatch-algo}

Once a proposition is classified, the \texttt{ScalarEncoder} in
\texttt{jugeo.encodings} translates the \Python{} expression into an
SMT-LIB assertion, and the \texttt{SolverRouter} dispatches it to Z3
with the appropriate logic declaration.

\begin{construction}[Encoding Pipeline]\label{con:encoding}
The encoding of a proposition $\prop$ at coordinate $c$ proceeds in
three stages:
\begin{enumerate}[nosep]
  \item \textbf{AST extraction.}
    Parse the \Python{} source at $c$ and extract the abstract syntax
    subtree corresponding to $\prop$.
  \item \textbf{Scalar encoding.}
    The \texttt{ScalarEncoder} translates \Python{} operations into
    SMT-LIB terms: arithmetic operators map to LIA/LRA operations,
    boolean connectives map to propositional logic, comparisons map
    to predicates, and function calls map to uninterpreted functions
    when no axiomatization is available.
  \item \textbf{Fragment-aware dispatch.}
    The classified fragment $F$ determines the \texttt{(set-logic F)}
    preamble, which restricts Z3 to the relevant theory solvers.
    The negation $\neg \prop$ is asserted, and a check-sat call
    determines satisfiability.
\end{enumerate}
\end{construction}

The dispatch logic respects the trust algebra:

\begin{lstlisting}[style=jugeo-python,
  caption={SMT dispatch with trust integration.},
  label={lst:dispatch}]
def dispatch_to_solver(
    prop: Proposition,
    coord: Coordinate,
    config: DescentConfiguration,
) -> EvidenceItem:
    smt_formula = ScalarEncoder.encode(prop, coord)
    fragment = classify_fragment(smt_formula)
    solver = z3.SolverFor(fragment.smt_logic_name)
    solver.set("timeout", int(config.overlap_check_timeout * 1000))
    solver.add(z3.Not(smt_formula.to_z3()))
    result = solver.check()
    if result == z3.unsat:
        return EvidenceItem(
            kind=EvidenceItemKind.SOLVER_PROOF,
            payload={"model": "refutation", "fragment": fragment.name},
            trust_level=TrustLevel.SOLVER_DISCHARGED,
            channel="z3",
        )
    elif result == z3.sat:
        model = solver.model()
        return EvidenceItem(
            kind=EvidenceItemKind.RUNTIME_WITNESS,
            payload={"counterexample": str(model)},
            trust_level=TrustLevel.CONTRADICTED,
            channel="z3",
        )
    else:  # unknown / timeout
        return EvidenceItem(
            kind=EvidenceItemKind.ORACLE_PROPOSAL,
            payload={"status": "timeout", "fragment": fragment.name},
            trust_level=TrustLevel.UNVERIFIED,
            channel="z3",
        )
\end{lstlisting}

\subsection{Encoding Theorems}\label{sec:encoding-thms}

\begin{theorem}[Encoding Soundness]\label{thm:encoding-sound}
Let $\prop$ be a proposition at coordinate $c$ in program $P$, and let
$E : \prop \mapsto \hat{\prop}$ be the scalar encoding into SMT-LIB
fragment $F$.  If the SMT solver returns $\mathsf{UNSAT}$ for
$\neg \hat{\prop}$ under encoding $E$, then $\prop$ holds in the
original \Python{} semantics with trust level $\Tsolver$.

Formally: $\mathsf{UNSAT}(\neg \hat{\prop})$ implies
$\llbracket \prop \rrbracket_P = \mathsf{true}$, provided $E$ is a
\emph{faithful encoding}---\ie it preserves and reflects satisfiability.
\end{theorem}

\begin{proof}
We establish faithfulness in two steps.

\emph{Preservation.}  The \texttt{ScalarEncoder} maps each \Python{}
operator to its SMT-LIB counterpart following the standard
interpretation: \texttt{+} to \texttt{bvadd} or \texttt{+} depending
on the sort, \texttt{==} to \texttt{=}, \texttt{and} to
\texttt{and}, and so forth.  Variables are declared with sorts
matching their \Python{} type annotations (or inferred types).
By construction, any assignment $\sigma$ satisfying $\prop$ in
\Python{} semantics also satisfies $\hat{\prop}$ in the SMT theory,
so $\mathsf{SAT}(\hat{\prop}) \Rightarrow \mathsf{SAT}(\prop)$.
Contrapositively,
$\mathsf{UNSAT}(\prop) \Rightarrow \mathsf{UNSAT}(\hat{\prop})$.

\emph{Reflection.}  Conversely, every model of $\hat{\prop}$ in the
SMT theory corresponds to a valid \Python{} execution, because the
encoding does not introduce fresh satisfying assignments: all SMT
variables correspond to program variables, and all constraints are
direct translations.  Thus
$\mathsf{UNSAT}(\hat{\prop}) \Rightarrow \mathsf{UNSAT}(\prop)$,
which yields
$\mathsf{UNSAT}(\neg\hat{\prop}) \Rightarrow
 \mathsf{UNSAT}(\neg\prop) \Rightarrow
 \llbracket \prop \rrbracket_P = \mathsf{true}$.

The trust level $\Tsolver$ is justified because the evidence was
produced by a sound decision procedure operating within a decidable
fragment $F$.
\end{proof}

\begin{theorem}[Fragment Optimality]\label{thm:frag-optimal}
Dispatching to the most specific fragment $F \sqsubseteq F'$ never
increases solving time and can decrease it by up to an order of
magnitude.
\end{theorem}

\begin{proof}[Proof sketch]
When Z3 receives a \texttt{(set-logic $F$)} declaration, it
activates only the theory solvers needed for $F$.  If $F$ is
\QFLIA{}, only the simplex engine is invoked; if $F$ is the less
specific $\texttt{QF\_AUFLIA}$, the array solver is additionally
activated even though the formula contains no array symbols.  The
overhead of unused theory solvers is measurable: in our benchmarks,
dispatching a pure \QFLIA{} formula as $\texttt{QF\_AUFLIA}$ incurs
a 3--10$\times$ slowdown due to unnecessary array-theory
propagation.  Since classification adds only $O(|\varphi|)$ overhead,
the net effect is always non-negative.
\end{proof}


%% ═══════════════════════════════════════════════════════════════════════
\section{The CLI and API}\label{sec:cli-api}
%% ═══════════════════════════════════════════════════════════════════════

\JuGeo{} exposes the entire geometric verification pipeline through
both a command-line interface (CLI) and a \Python{} API.  The CLI
provides high-level entry points suitable for continuous integration
and interactive development; the API provides fine-grained control over
sites, judgments, and descent for programmatic use.

\subsection{CLI Interface}\label{sec:cli}

The \texttt{jugeo} command dispatches to six primary subcommands, each
targeting a distinct stage of the pipeline:

\begin{lstlisting}[style=jugeo-output,
  caption={\JuGeo{} CLI subcommands.},
  label={lst:cli}]
jugeo prove FILE.py          # Full descent verification
jugeo encode FILE.py         # SMT encoding only
jugeo bugs FILE.py           # Bug detection via H^1
jugeo repair FILE.py         # Obstruction-guided repair
jugeo equiv FILE1 FILE2      # Equivalence checking
jugeo synth SPEC.py          # Cover-guided synthesis
\end{lstlisting}

\paragraph{\texttt{jugeo prove}.}
The primary verification command.  Given a \Python{} source file,
\texttt{prove} executes the full descent pipeline
(\Cref{sec:pipeline}): parses the AST, constructs the semantic site,
computes covers, encodes propositions to SMT, dispatches to Z3,
checks \v{C}ech cohomology, and emits a certificate if
$\Hcoh{1} = 0$.

\paragraph{\texttt{jugeo encode}.}
Runs only the encoding phase, producing SMT-LIB assertions without
solving.  Useful for debugging encodings, feeding to external
solvers, or inspecting fragment classification.

\paragraph{\texttt{jugeo bugs}.}
Computes $\Hcoh{1}(\Site(P), \J)$ and reports each non-trivial
cohomology class as a potential bug.  Each report includes the
obstruction coordinate, the violated overlap condition, the
cohomological class, severity, and repair hints.

\paragraph{\texttt{jugeo repair}.}
Extends \texttt{bugs} with the repair algorithm of
\Cref{sec:repair}: for each $[\alpha] \in \Hcoh{1}$, the tool
classifies the obstruction, localizes it, proposes a code patch,
and re-verifies.  The output includes the original obstruction, the
proposed patch, and the post-repair trust level.

\paragraph{\texttt{jugeo equiv}.}
Checks behavioral equivalence of two \Python{} files by constructing
a joint site with a shared interface cover, verifying each
implementation against the shared specification, and gluing.

\paragraph{\texttt{jugeo synth}.}
Given a specification file, synthesizes a \Python{} implementation
guided by the covering structure: the specification is decomposed
into local constraints via the cover, each constraint is solved
independently (by SMT or enumerative search), and the local
solutions are glued into a complete program.

\subsection{Python API}\label{sec:api}

The deep API exposes all layers of the geometric framework.  We
illustrate the core workflow: building a site, running descent, and
constructing judgments.

\begin{lstlisting}[style=jugeo-python,
  caption={Site construction via the \Python{} API.},
  label={lst:api-site}]
from jugeo.geometry.site import (
    SiteBuilder, Coordinate, CoordinateKind,
    Morphism, MorphismKind,
)
from jugeo.geometry.covers import Cover, CoverBuilder
from jugeo.geometry.descent import (
    DescentEngine, DescentConfiguration,
    DescentStrategy,
)

# Build a semantic site from program structure
site = (SiteBuilder("my_program")
    .add_coordinate(Coordinate(
        components=("my_program", "main"),
        kind=CoordinateKind.FUNCTION))
    .add_coordinate(Coordinate(
        components=("my_program", "helper"),
        kind=CoordinateKind.FUNCTION))
    .add_morphism(Morphism(
        source=Coordinate(("my_program", "main"),
                          CoordinateKind.FUNCTION),
        target=Coordinate(("my_program", "helper"),
                          CoordinateKind.FUNCTION),
        kind=MorphismKind.RESTRICTION,
        label="main calls helper"))
    .build())
\end{lstlisting}

\begin{lstlisting}[style=jugeo-python,
  caption={Running descent and inspecting results.},
  label={lst:api-descent}]
# Configure and run descent
engine = DescentEngine(
    configuration=DescentConfiguration(
        strategy=DescentStrategy.ITERATIVE,
        depth_limit=5,
        copilot_assisted_repair=True,
    )
)

# attempt_descent returns a DescentResult:
#   either GlobalSection (success) or
#   DescentObstruction (failure with H^1 class)
result = engine.attempt_descent(
    cover=site.covering_families()[0],
    sections=local_sections,
)

if result.is_success:
    section = result.unwrap_section()
    print(f"Verified: trust = {section.trust_floor}")
    print(f"Certificate: {section.certificate}")
else:
    obs = result.unwrap_obstruction()
    print(f"Obstruction: {obs.cohomology_class.summary()}")
    print(f"Repair frontier: {obs.repair_frontier.summary()}")
\end{lstlisting}

\begin{lstlisting}[style=jugeo-python,
  caption={Judgment construction via the fluent builder.},
  label={lst:api-judgment}]
from jugeo.judgments.judgment_terms import (
    Judgment, JudgmentBuilder,
    Proposition, PropositionKind,
    TrustLevel, EvidenceItem, EvidenceItemKind,
)

# Build a judgment (the 8-tuple)
j = (JudgmentBuilder()
    .at(Coordinate(("my_program", "main"),
                   CoordinateKind.FUNCTION))
    .claiming(Proposition(
        kind=PropositionKind.BEHAVIORAL,
        formula="returns sorted list",
    ))
    .of_type_named("FunctionContract")
    .with_trust_level(TrustLevel.SOLVER_DISCHARGED)
    .with_evidence(EvidenceItem(
        kind=EvidenceItemKind.SOLVER_PROOF,
        payload={"solver": "z3", "fragment": "QF_LIA"},
        trust_level=TrustLevel.SOLVER_DISCHARGED,
        channel="z3",
    ))
    .build())

# Judgment operations
restricted = j.restrict_to(
    Coordinate(("my_program", "main", "loop_body"),
               CoordinateKind.REGION))
transported = j.transport_along(equivalence_morphism)
\end{lstlisting}

\subsection{Pipeline Architecture}\label{sec:pipeline}

The full verification pipeline composes the components of
\Cref{sec:moves,sec:dispatch,sec:cli} into a nine-stage process.
\Cref{fig:pipeline} shows the data flow.

\begin{figure}[htbp]
\centering
\begin{tikzcd}[
  row sep=2em, column sep=2em,
  cells={font=\small},
  every arrow/.append style={-stealth}
]
  \text{\Python{} Source}
    \arrow[r]
  & \text{AST} \to \Site_P
    \arrow[r]
  & \Cov \to \text{SMT}
    \arrow[d] \\
  \text{Certificate}
  & \Hcoh{1}
    \arrow[l, "{=\,0}"]
    \arrow[r, "{\neq\,0}"']
  & \text{Z3 Verdicts}
    \arrow[l] \\
  & \REPAIR
    \arrow[uur, bend right=25, dashed, "\text{iterate}"']
  &
\end{tikzcd}
\caption{The \JuGeo{} verification pipeline.  Solid arrows show the
  primary data flow; the dashed arrow shows the repair loop.}
\label{fig:pipeline}
\end{figure}

The nine stages are:

\begin{enumerate}[label=\textbf{Stage \arabic*.},leftmargin=*]
  \item \textbf{Parse \Python{} AST.}
    The input file is parsed using \Python's built-in \texttt{ast}
    module.  Each function, class, and control-flow branch is
    identified.

  \item \textbf{Extract coordinates and morphisms.}
    Functions become coordinates of kind $\mathsf{FUNCTION}$; classes
    become $\mathsf{MODULE}$; loop bodies and conditional branches
    become $\mathsf{REGION}$.  Call edges, scoping relations, and
    refinement arrows become morphisms of kind $\mathsf{RESTRICTION}$,
    $\mathsf{INCLUSION}$, $\mathsf{TRANSPORT}$, or
    $\mathsf{REFINEMENT}$.

  \item \textbf{Build semantic site $\Site(P)$.}
    The \texttt{SiteBuilder} assembles the category
    $\mathcal{C}_P$ and equips it with the canonical Grothendieck
    topology (satisfying axioms T1--T3 from \Cref{sec:sites}).

  \item \textbf{Compute covering families.}
    The \texttt{CoverGenerator} produces covers via four strategies:
    \texttt{from\_scope\_tree()} for scope covers,
    \texttt{from\_call\_graph()} for call-graph covers,
    \texttt{from\_module\_structure()} for module covers, and
    \texttt{from\_class\_hierarchy()} for class-hierarchy covers.
    Each cover is validated against the Grothendieck axioms via
    \texttt{CoverDiagnostics.check\_covering\_axiom()}.

  \item \textbf{Extract local sections.}
    For each patch $c_i$ in a covering family, the pipeline
    extracts a local section: the set of propositions (preconditions,
    postconditions, invariants, type constraints) that are asserted at
    $c_i$.  Each section is a \texttt{LocalSection} carrying
    the judgment data, evidence bundle, and trust level.

  \item \textbf{Encode to SMT.}
    The \texttt{ScalarEncoder} translates each proposition into an
    SMT-LIB formula.  Fragment classification (\Cref{sec:fragment-class})
    determines the minimal decidable fragment.

  \item \textbf{Dispatch to Z3 and collect verdicts.}
    Each encoded formula is dispatched to Z3 with the optimal
    \texttt{(set-logic $F$)} declaration.  The solver returns
    $\mathsf{UNSAT}$ (valid), $\mathsf{SAT}$ (counterexample), or
    $\mathsf{UNKNOWN}$ (timeout).  Each verdict becomes an
    \texttt{EvidenceItem} with the appropriate trust level
    (\Cref{lst:dispatch}).

  \item \textbf{Check \v{C}ech cohomology.}
    The \texttt{DescentEngine} checks the overlap conditions
    and computes $\Hcoh{1}(\Cech(\mathcal{U}), \J)$.
    If $\Hcoh{1} = 0$: the $\GLUE$ move produces a
    \texttt{GlobalSection} and a verification certificate.
    If $\Hcoh{1} \neq 0$: each non-trivial class $[\alpha]$
    is recorded as a \texttt{DescentObstruction} with its
    \texttt{CohomologyClass} and \texttt{RepairFrontier}.

  \item \textbf{Emit certificate or repair.}
    On success, the pipeline emits a machine-checkable certificate
    recording the site, the cover, the local verdicts, the overlap
    evidence, and the glued global section.  On failure, the pipeline
    either reports the obstructions (in \texttt{bugs} mode) or
    invokes the repair algorithm (\Cref{sec:repair}) and iterates.
\end{enumerate}

\begin{proposition}[Pipeline Correctness]\label{prop:pipeline}
The nine-stage pipeline is a composition of semantic moves:
\[
  \mathsf{pipeline} \;=\;
  \GLUE \;\circ\;
  \bigl(\VERIFY_{c_1} \otimes \cdots \otimes \VERIFY_{c_n}\bigr)
  \;\circ\;
  \bigl(\RESTRICT_{c_1} \otimes \cdots \otimes \RESTRICT_{c_n}\bigr)
\]
where $\{c_1, \ldots, c_n\}$ is the covering family computed in
Stage~4.  By \Cref{thm:move-soundness}, if the pipeline succeeds,
the resulting certificate is sound.
\end{proposition}


%% ═══════════════════════════════════════════════════════════════════════
\section{Program Repair via Obstruction Surgery}\label{sec:repair}
%% ═══════════════════════════════════════════════════════════════════════

When $\Hcoh{1}(\Site(P), \J) \neq 0$, each non-trivial cohomology
class $[\alpha]$ represents an obstruction to gluing: a concrete
failure in the local-to-global transfer.  Rather than merely reporting
these failures, \JuGeo{} attempts to \emph{repair} them---to modify
the program $P$ so that the obstructions vanish and descent succeeds.
We call this process \textbf{obstruction surgery}, by analogy with
the topological operation of ``killing'' homotopy classes via surgery
on manifolds.

\subsection{The Repair Algorithm}\label{sec:repair-algo}

The repair algorithm operates on the set
$\mathrm{Obs}(\Site(P)) = \{[\alpha_1], \ldots, [\alpha_k]\}$ of
non-trivial classes in $\Hcoh{1}$.  It proceeds in five steps.

\begin{construction}[Obstruction Surgery]\label{con:repair}
Given a program $P$ with $\Hcoh{1}(\Site(P), \J) \neq 0$:

\begin{enumerate}[label=\textbf{Step \arabic*.},leftmargin=*]
  \item \textbf{Classify.}
    For each $[\alpha] \in \Hcoh{1}$, determine the obstruction kind
    from its cocycle data.  The classification follows the
    \texttt{Obstruction} dataclass:
    \begin{itemize}[nosep]
      \item \textsc{Type\_Mismatch}: the overlap condition fails
            because types disagree across patches (\eg
            \texttt{int} vs.\ \texttt{float} at a shared interface).
      \item \textsc{Missing\_Case}: a branch is absent from one
            patch, causing the overlap predicate to be undefined.
      \item \textsc{Invariant\_Violation}: a loop invariant
            established in one patch is violated in an adjacent patch.
      \item \textsc{Contract\_Breach}: a precondition or
            postcondition is not satisfied at the overlap.
      \item \textsc{Aliasing\_Conflict}: two patches share mutable
            state, and their modifications are incompatible.
    \end{itemize}

  \item \textbf{Localize.}
    Find the \emph{minimal coordinate} where $[\alpha]$ is
    supported.  The cocycle $\alpha$ is a function on pairs of
    overlapping patches; its support is the set of overlap
    coordinates where $\alpha$ is non-trivial.  The minimal support
    coordinate is the common ancestor
    $c_\alpha = \bigwedge_{(i,j) \in \mathrm{supp}(\alpha)}
     (c_i \wedge c_j)$
    in the coordinate lattice.

  \item \textbf{Suggest patch.}
    Generate a code modification $\sigma : P \to P'$ at $c_\alpha$
    that ``kills'' $[\alpha]$---\ie makes $\alpha$ a coboundary in
    the modified \v{C}ech complex.  The patch is constructed from the
    repair hints stored in the obstruction record
    (\texttt{Obstruction.repair\_hints}) and the localized coordinate.
    For $\textsc{Type\_Mismatch}$ obstructions, $\sigma$ inserts a
    type coercion; for $\textsc{Missing\_Case}$, $\sigma$ adds a
    branch; for $\textsc{Contract\_Breach}$, $\sigma$ strengthens or
    weakens the contract.

  \item \textbf{Verify.}
    Re-run the descent pipeline on $P'$.  Specifically, recompute
    the local sections affected by $\sigma$ (those at or below
    $c_\alpha$), re-encode, re-dispatch to Z3, and re-check the
    overlap conditions.  The repair frontier
    (\texttt{RepairFrontier}) ensures that only the affected portion
    of the site is re-verified, keeping the cost proportional to the
    support of $[\alpha]$ rather than the entire program.

  \item \textbf{Iterate.}
    If $\Hcoh{1}(\Site(P'), \J) = 0$, the repair succeeds and the
    pipeline proceeds to gluing.  If not, return to Step~1 with
    $P \leftarrow P'$ and the reduced obstruction set.
\end{enumerate}
\end{construction}

The following diagram illustrates the surgical process:

\begin{center}
\begin{tikzcd}[column sep=3em, row sep=2em]
  P \arrow[r, "\text{descent}"]
    \arrow[d, "\sigma_1"']
  & \Hcoh{1} \neq 0
    \arrow[d, "\text{classify}"] \\
  P' \arrow[r, "\text{descent}"]
    \arrow[d, dashed, "\sigma_2"']
  & {[\alpha_1] = 0,\; [\alpha_2] \neq 0}
    \arrow[d, "\text{classify}"] \\
  P'' \arrow[r, "\text{descent}"]
  & \Hcoh{1} = 0 \;\checkmark
\end{tikzcd}
\end{center}

\begin{example}[Repair of a type mismatch]\label{ex:type-repair}
Consider two functions sharing an interface:

\begin{lstlisting}[style=jugeo-python,caption={Buggy program with type mismatch.}]
def producer() -> int:
    return 42

def consumer(x: float) -> float:
    return x * 1.5
\end{lstlisting}

\noindent
The semantic site has two patches (one per function) overlapping at the
shared variable \texttt{x}.  The overlap condition requires
$\res_{\text{producer}}(\mathtt{type}(x)) =
 \res_{\text{consumer}}(\mathtt{type}(x))$, but
$\mathtt{int} \neq \mathtt{float}$, yielding a non-trivial
$[\alpha] \in \Hcoh{1}$ classified as $\textsc{Type\_Mismatch}$.

The repair algorithm localizes $[\alpha]$ to the call site, suggests
$\sigma$: insert \texttt{float(producer())} at the call, and
re-verifies.  After the patch, the overlap condition is satisfied
and $\Hcoh{1} = 0$.
\end{example}

\subsection{Repair Soundness}\label{sec:repair-soundness}

We now establish that the repair algorithm terminates and produces
a correct result.

\begin{definition}[Cocycle Norm]\label{def:cocycle-norm}
For a cohomology class $[\alpha] \in \Hcoh{1}$, define its
\emph{rank} $\mathrm{rk}([\alpha])$ as the number of overlap pairs
$(i,j)$ where $\alpha_{ij}$ is non-trivial.  Define the
\emph{total obstruction norm} as
$\lVert \mathrm{Obs}(\Site(P)) \rVert =
 \sum_{[\alpha]} \mathrm{rk}([\alpha])$.
\end{definition}

\begin{lemma}[Norm Reduction]\label{lem:norm-reduce}
Each successful repair step reduces the total obstruction norm by at
least one.  That is, if $\sigma$ kills $[\alpha]$ and does not
introduce new obstructions, then
$\lVert \mathrm{Obs}(\Site(P')) \rVert <
 \lVert \mathrm{Obs}(\Site(P)) \rVert$.
\end{lemma}

\begin{proof}
A repair $\sigma$ that kills $[\alpha]$ makes $\alpha$ a coboundary:
there exists a 0-cochain $\beta$ with $\alpha = \delta\beta$.  The
class $[\alpha]$ vanishes from $\Hcoh{1}(\Site(P'), \J)$, removing
$\mathrm{rk}([\alpha]) \geq 1$ non-trivial overlap pairs from the
total norm.  The assumption that $\sigma$ does not introduce new
obstructions (verified in Step~4 of \Cref{con:repair}) ensures
that no new classes appear.
\end{proof}

\begin{theorem}[Repair Termination]\label{thm:repair-term}
Under the assumption that each repair step does not introduce new
obstructions (verified by re-running descent), the repair algorithm
terminates in at most $\lvert \mathrm{Obs}(\Site(P)) \rvert$ steps.
\end{theorem}

\begin{proof}
The total obstruction norm
$\lVert \mathrm{Obs}(\Site(P)) \rVert \in \mathbb{N}$ is a
well-founded measure.  By \Cref{lem:norm-reduce}, each iteration
strictly decreases the norm.  Since $\lVert \cdot \rVert \geq 0$,
the algorithm terminates after at most
$\lVert \mathrm{Obs}(\Site(P_0)) \rVert$ iterations, where $P_0$ is
the original program.

In the worst case, each iteration kills exactly one obstruction class
of rank~1, so the bound is $|\mathrm{Obs}(\Site(P_0))|$.  In
practice, a single repair often kills multiple obstruction classes
simultaneously (when they share a support coordinate), and the
algorithm terminates faster.
\end{proof}

\begin{remark}[Non-introduction of obstructions]\label{rem:non-intro}
The assumption that repairs do not introduce new obstructions is
enforced by Step~4 (re-verification).  If a repair \emph{does}
introduce a new class $[\alpha']$, then $\lVert \mathrm{Obs} \rVert$
may not decrease, and the algorithm falls back to reporting the new
obstruction rather than looping.  In the worst case, the algorithm
halts after a configurable iteration budget (default: $2 \times
|\mathrm{Obs}(\Site(P_0))|$) and reports all remaining obstructions.
\end{remark}

\begin{theorem}[Repair Correctness]\label{thm:repair-correct}
If the repair algorithm terminates with $\Hcoh{1}(\Site(P'), \J) = 0$,
then the descent theorem applies: the local sections glue to a global
section, and the pipeline emits a valid certificate for $P'$.
\end{theorem}

\begin{proof}
By the descent theorem (\Cref{sec:descent}), $\Hcoh{1} = 0$ implies
that compatible local sections glue uniquely.  The re-verification in
Step~4 ensures that all local sections of $P'$ are valid (each has
been re-dispatched to Z3 and confirmed).  The overlap conditions hold
(they are exactly the cocycle condition whose triviality is asserted
by $\Hcoh{1} = 0$).  Therefore the $\GLUE$ move succeeds, producing
a \texttt{GlobalSection} with a valid certificate.  By
\Cref{thm:move-soundness}, the certificate is sound.
\end{proof}

\begin{corollary}[End-to-End Pipeline Soundness]\label{cor:e2e}
If the \JuGeo{} pipeline (with or without repair) terminates
successfully and emits a certificate, then the verified propositions
hold in the original (or repaired) \Python{} program.
\end{corollary}

\begin{proof}
Combine \Cref{thm:encoding-sound} (encoding faithfulness),
\Cref{thm:move-soundness} (move soundness),
\Cref{prop:pipeline} (pipeline as move composition), and
\Cref{thm:repair-correct} (repair correctness).  Each stage preserves
soundness, so the composition is sound.
\end{proof}
